\documentclass[aspectratio=169]{beamer}
\usepackage[utf8]{inputenc}
\usepackage[utf8]{vietnam}
\usepackage{hyperref}
\usepackage{tcolorbox}
\usepackage{booktabs}

% Theme selection
\usetheme{Madrid}
\usecolortheme{default}

\title{Week 2 Progress: LLM Security \& Prompt Injection Defense}
\subtitle{Quantifying Vulnerabilities and Implementing Defense Layers}
\author{Your Name / Group}
\date{\today}

\begin{document}

% Title Slide
\begin{frame}
    \titlepage
\end{frame}

% Section 1: Scaled Attack Results (First Slide hook)
\begin{frame}{Scaled Attack Results: The Threat is Real}
    \textbf{Executive Summary:} We expanded our attack campaign significantly from the pilot phase to quantify the true vulnerability of SOTA models. \\~\\
    
    \begin{columns}[T]
        \begin{column}{0.55\textwidth}
            \textbf{Testing Parameters:}
            \begin{itemize}
                \item \textbf{Models Evaluated:} GPT OSS 120B and Owl Alpha.
                \item \textbf{Corpus:} Scaled up to hundreds of diverse medical documents.
                \item \textbf{Variables:} Tested across 3 Difficulty Tiers (Easy, Medium, Hard) and 3 Injection Positions (Beginning, Middle, End).
            \end{itemize}
        \end{column}
        
        \begin{column}{0.45\textwidth}
            \begin{tcolorbox}[colback=red!5,colframe=red!75!black,title=Critical Vulnerability]
            \centering
            Overall Attack Success Rate:\\
            \vspace{0.2cm}
            \huge\textbf{74.1\%}
            \end{tcolorbox}
        \end{column}
    \end{columns}
    
    \vspace{0.3cm}
    \textbf{Conclusion:} Relying on system instructions is fundamentally broken. We must implement active defense layers.
\end{frame}

% Section 2: Stats Breakdown 1
\begin{frame}{Results Breakdown: Models \& Difficulty}
    \begin{columns}[T]
        \begin{column}{0.48\textwidth}
            \begin{block}{ASR by Target Model}
                \centering
                \renewcommand{\arraystretch}{1.3}
                \begin{tabular}{lc}
                    \toprule
                    \textbf{Model} & \textbf{Overall ASR} \\
                    \midrule
                    GPT OSS 120B & \textbf{86.7\%} \\
                    Owl Alpha & 82.2\% \\
                    Gemini 3.1 Flash Lite & 14.7\% \\
                    \bottomrule
                \end{tabular}
            \end{block}
            \vspace{0.2cm}
            \small\textit{Open-weights models were critically compromised in over 80\% of all attacks. Proprietary models (Gemini) exhibited tighter alignment but were not immune.}
        \end{column}
        
        \begin{column}{0.48\textwidth}
            \begin{block}{ASR by Template Difficulty}
                \centering
                \renewcommand{\arraystretch}{1.3}
                \begin{tabular}{lc}
                    \toprule
                    \textbf{Difficulty Tier} & \textbf{Peak ASR} \\
                    \midrule
                    Medium (Structural Deception) & \textbf{86.8\%} \\
                    Easy (Direct Override) & 63.1\% \\
                    Hard (Stealth / Obfuscation) & 41.3\% \\
                    \bottomrule
                \end{tabular}
            \end{block}
            \vspace{0.2cm}
            \small\textit{Structural deception (Medium) proved to be the most lethal vector, completely confusing the LLM's understanding of system vs user boundaries.}
        \end{column}
    \end{columns}
\end{frame}

% Section 3: Stats Breakdown 2
\begin{frame}{PII Extraction Deep-Dive: The Position Paradox}
    \textbf{Focus:} Analyzing ONLY the \textit{PII Extraction} attack category across the full historical dataset (680+ attacks). \\~\\
    
    \begin{block}{PII Leakage by Payload Position}
        \centering
        \renewcommand{\arraystretch}{1.3}
        \begin{tabular}{lc|c}
            \toprule
            \textbf{Position} & \textbf{PII Leak Rate} & \textbf{Observation} \\
            \midrule
            \textbf{Beginning} & \textbf{88.7\%} & High compliance (Primacy effect). \\
            \textbf{End} & \textbf{92.8\%} & Highest compliance (Recency effect). \\
            \textbf{Middle} & \textbf{87.3\%} & Sustained high compliance. \\
            \bottomrule
        \end{tabular}
    \end{block}
    
    \vspace{0.3cm}
    \textbf{Critical Finding:} While general prompt injections suffer from the "Lost in the Middle" phenomenon, \textbf{PII Extraction completely bypasses it}. Models are so eager to comply with data extraction requests that they will hunt for and leak PII regardless of where the malicious payload is hidden in the document.
\end{frame}



% Section 10: Privacy Filter Architecture
\begin{frame}{Deep Dive: OPF Architecture}
    \textbf{The Model:} We integrated the official OpenAI Privacy Filter (\texttt{opf}), a highly optimized, open-weight neural network designed specifically for PII redaction. \\~\\
    
    \begin{columns}[T]
        \begin{column}{0.48\textwidth}
            \begin{block}{Architecture}
                \begin{itemize}
                    \item \textbf{Type:} Lightweight, \textbf{bidirectional token-classification model}.
                    \item \textbf{Scale:} Designed to be extremely fast and efficient, allowing it to run inference over large documents with minimal latency overhead.
                \end{itemize}
            \end{block}
        \end{column}
        
        \begin{column}{0.48\textwidth}
            \begin{block}{Training Methodology}
                \begin{itemize}
                    \item \textbf{Multilingual Support:} Trained on massive corpuses of diverse, multilingual text.
                    \item \textbf{Synthetic Augmentation:} Heavily augmented with algorithmically generated synthetic PII to prevent the model from overfitting.
                \end{itemize}
            \end{block}
        \end{column}
    \end{columns}
\end{frame}

% Section 11: Privacy Filter Output
\begin{frame}{Deep Dive: Output Shape \& Decoding}
    \textbf{The Pipeline:} The model frames PII detection as a dense token-classification task across 8 distinct privacy categories (e.g., \textit{account\_number, private\_address, private\_person}). \\~\\
    
    \begin{block}{BIOES Tagging \& Logit Output}
        \begin{itemize}
            \item \textbf{Class Expansion:} Each of the 8 categories is expanded into boundary-tagged classes using the \textbf{BIOES} format (\textbf{B}egin, \textbf{I}nside, \textbf{E}nd, \textbf{S}ingle).
            \item \textbf{Output Space:} 1 Background Class ($O$) + (8 categories $\times$ 4 boundary tags) = \textbf{33 total output classes}.
            \item \textbf{Logit Shape:} For a sequence of length $T$, the output head emits a tensor of shape \textbf{$[T, 33]$}. For batched inference, this becomes \textbf{$[B, T, 33]$}.
        \end{itemize}
    \end{block}
    
    \vspace{0.2cm}
    \textbf{Constrained Sequence Decoding:} At inference time, these 33 raw per-token logits are not just greedily selected. They are decoded into coherent, non-overlapping span labels using a constrained sequence decoding algorithm, mathematically guaranteeing valid BIOES transitions.
\end{frame}

% Section 12: DPO Architecture
\begin{frame}{Defense Layer 2: Behavioral Alignment via DPO}
    \textbf{The Concept:} Teach the AI to actively recognize and refuse malicious instructions by penalizing compliance and rewarding safe summarization. \\~\\
    
    \begin{block}{Model Architecture \& LoRA Finetuning}
        \begin{itemize}
            \item \textbf{Base Model:} Selected \texttt{Qwen/Qwen2.5-1.5B-Instruct} as the foundation. At 1.5 billion parameters, it is highly capable yet small enough for rapid experimentation.
            \item \textbf{Low-Rank Adaptation (LoRA):} Instead of performing a computationally expensive full-parameter update, we utilized LoRA combined with 4-bit quantization (\texttt{bitsandbytes}).
            \item \textbf{Efficiency:} LoRA only trains a tiny fraction of parameter matrices (the adapter), reducing VRAM requirements to fit comfortably on Kaggle T4 GPUs while preventing catastrophic forgetting.
        \end{itemize}
    \end{block}
\end{frame}

% Section 13: DPO Data Selection
\begin{frame}{Deep Dive: DPO Preference Data Selection}
    \textbf{The Challenge:} Direct Preference Optimization requires pairs of responses—one \textit{Chosen} (good) and one \textit{Rejected} (bad). \\~\\
    
    \begin{block}{Data Mining Strategy}
        We synthetically generated a high-quality dataset by mining our own Attack Master Logs:
        \begin{itemize}
            \item \textbf{The Input:} A document heavily poisoned with a prompt injection payload.
            \item \textbf{The Rejected Response (Penalty):} We filtered the logs for \textit{successful} attacks. The output where the model leaked PII or hallucinated was set as the mathematically penalized response.
            \item \textbf{The Chosen Response (Reward):} We paired it with the \textit{clean baseline summary} (where the model successfully summarized the text and ignored the injection).
        \end{itemize}
    \end{block}
    
    \vspace{0.2cm}
    \small\textit{This highly specific data selection forces the model to unlearn its obedience to injected commands, aligning its internal probabilities strictly toward safe summarization.}
\end{frame}

% Section 13: Defense Evaluation Results
\begin{frame}{Defense Evaluation Results}
    \textbf{The Benchmark:} Evaluated 50 highly aggressive PII Extraction attacks specifically against the unprotected Base Model versus the Privacy Filter defense layer. \\~\\
    
    \begin{block}{Defense Performance (ASR)}
        \centering
        \renewcommand{\arraystretch}{1.3}
        \begin{tabular}{lc|c}
            \toprule
            \textbf{System Configuration} & \textbf{Attack Success Rate} & \textbf{Status} \\
            \midrule
            \textbf{Base Model (No Defense)} & \textbf{70.0\%} & Highly Vulnerable (35/50 leaked) \\
            \textbf{Baseline Privacy Filter} & \textbf{8.0\%} & \textbf{Massive Mitigation} (Only 4/50 leaked) \\
            \textbf{DPO-Aligned Model} & \textit{Training...} & Pending final epoch completion \\
            \bottomrule
        \end{tabular}
    \end{block}
    
    \vspace{0.3cm}
    \textbf{Conclusion:} 
    \begin{itemize}
        \item The Privacy Filter successfully intercepted and neutralized \textbf{88.5\%} of lethal attacks that would have breached the base model.
        \item \textbf{Next Step:} Complete DPO training to determine if behavioral alignment can achieve absolute 0\% ASR.
    \end{itemize}
\end{frame}

% Section 14: Future Direction (AdvGRPO)
\begin{frame}{Future Work: Adaptive Red Teaming via GRPO}
    \textbf{Research Expansion:} Exploring state-of-the-art co-training methodologies to evolve our defense layers beyond static DPO alignment. \\~\\
    
    \begin{block}{Paper Highlight: \textit{Learning to Attack and Defend} (AdvGRPO)}
        \begin{itemize}
            \item \textbf{The Core Concept:} Instead of training a defense against a static dataset of attacks, we simultaneously co-train an \textbf{Attacker Model} and a \textbf{Defender Model} in tandem using Reinforcement Learning (GRPO).
            \item \textbf{The Mechanism:} The Attacker is rewarded for successfully bypassing the Defender. The Defender is rewarded for safely neutralizing the Attacker while maintaining utility.
            \item \textbf{Why it Matters:} This creates an "arms race" that automatically discovers novel zero-day prompt injection techniques and patches them dynamically.
        \end{itemize}
    \end{block}
    
    \vspace{0.3cm}
    \textbf{Implementation Goal:} Upgrade our current static DPO pipeline to an active AdvGRPO co-training loop, producing a defense layer that autonomously adapts to evolving prompt injection strategies.
\end{frame}

\end{document}
